\documentclass[aspectratio=169,11pt]{beamer}
\usetheme{default}
\setbeamertemplate{navigation symbols}{}
\setbeamertemplate{footline}{%
  \hfill\insertframenumber/\inserttotalframenumber\hspace*{4pt}\vspace{2pt}%
}
\usepackage{lmodern}
\usepackage[most]{tcolorbox}
\usepackage{listings}
\usepackage{tikz}
\usetikzlibrary{arrows.meta,positioning,shapes.geometric,fit,backgrounds}
\usepackage{booktabs}
\usepackage{hyperref}

\definecolor{lcgreen}{HTML}{2E7D32}
\definecolor{lcblue}{HTML}{1565C0}
\definecolor{lcorange}{HTML}{EF6C00}
\definecolor{lcred}{HTML}{C62828}
\definecolor{codebg}{HTML}{F5F5F5}
\definecolor{docgray}{gray}{0.55}

\newtcolorbox{conceptbox}[1][]{colback=yellow!20,colframe=yellow!60!black,title=#1,fonttitle=\bfseries}
\newtcolorbox{warningbox}[1][]{colback=red!15,colframe=red!60!black,title=#1,fonttitle=\bfseries}
\newtcolorbox{tipbox}[1][]{colback=green!10,colframe=green!50!black,title=#1,fonttitle=\bfseries}
\newtcolorbox{codebox}[1][]{colback=blue!10,colframe=blue!40!black,title=#1,fonttitle=\bfseries}

\lstset{
  basicstyle=\ttfamily\scriptsize,
  backgroundcolor=\color{codebg},
  breaklines=true,
  frame=single,
  rulecolor=\color{blue!40!black},
  keywordstyle=\color{lcblue}\bfseries,
  stringstyle=\color{lcgreen},
  commentstyle=\color{docgray}\itshape,
  language=Python,
  showstringspaces=false,
  tabsize=2,
  xleftmargin=2pt,
  xrightmargin=2pt,
}

\newcommand{\docref}[1]{{\scriptsize\color{docgray}[Docs: #1]}}

\title{LangChain v1.2.x Quick Reference}
\subtitle{Condensed Cheat Sheet • March 2026}
\author{}
\date{}

\begin{document}

% SLIDE 1: Title
\begin{frame}
  \titlepage
  \vspace{-1cm}
  \begin{center}
    {\small \textcolor{docgray}{A quick-lookup reference for the 20 most essential LangChain patterns}}
  \end{center}
\end{frame}

% SLIDE 2: Essential Imports
\begin{frame}[fragile]
  \frametitle{Essential Imports (Top 20)}

  \begin{columns}[T]
    \column{0.5\textwidth}
    \textbf{Models \& Chat}
    \begin{lstlisting}
from langchain_openai import ChatOpenAI
from langchain_anthropic import ChatAnthropic
from langchain_google_vertexai import ChatVertexAI
from langchain_ollama import ChatOllama
    \end{lstlisting}

    \textbf{Prompts}
    \begin{lstlisting}
from langchain_core.prompts import (
  ChatPromptTemplate,
  MessagesPlaceholder,
  PromptTemplate
)
    \end{lstlisting}

    \column{0.5\textwidth}
    \textbf{Chains \& Tools}
    \begin{lstlisting}
from langchain_core.tools import tool
from langchain.agents import (
  create_react_agent,
  AgentExecutor
)
    \end{lstlisting}

    \textbf{RAG \& Memory}
    \begin{lstlisting}
from langchain_text_splitters import (
  RecursiveCharacterTextSplitter
)
from langchain_core.runnables import (
  RunnableWithMessageHistory,
  RunnablePassthrough
)
from langchain_chroma import Chroma
    \end{lstlisting}
  \end{columns}
\end{frame}

% SLIDE 3: More Imports
\begin{frame}[fragile]
  \frametitle{Essential Imports (Continued)}

  \begin{columns}[T]
    \column{0.5\textwidth}
    \textbf{Outputs \& Parsing}
    \begin{lstlisting}
from langchain_core.output_parsers import (
  JsonOutputParser,
  StrOutputParser,
  PydanticOutputParser
)
from langchain.callbacks.tracers.langchain import (
  LangChainTracer
)
    \end{lstlisting}

    \column{0.5\textwidth}
    \textbf{LangGraph}
    \begin{lstlisting}
from langgraph.graph import StateGraph
from langgraph.prebuilt import (
  create_react_agent
)
    \end{lstlisting}

    \textbf{Documents \& Loaders}
    \begin{lstlisting}
from langchain_community.document_loaders import (
  PyPDFLoader,
  WebBaseLoader,
  CSVLoader
)
from langchain_core.documents import Document
    \end{lstlisting}
  \end{columns}
\end{frame}

% SLIDE 4: Chat Models
\begin{frame}[fragile]
  \frametitle{Chat Models — Quick Init}

  \begin{lstlisting}
# OpenAI
from langchain_openai import ChatOpenAI
llm = ChatOpenAI(model="gpt-4o-mini", temperature=0, api_key="...")

# Anthropic
from langchain_anthropic import ChatAnthropic
llm = ChatAnthropic(model="claude-3-5-sonnet-20241022", temperature=0)

# Google VertexAI
from langchain_google_vertexai import ChatVertexAI
llm = ChatVertexAI(model="gemini-2.0-flash", temperature=0)

# Ollama (Local)
from langchain_ollama import ChatOllama
llm = ChatOllama(model="llama2", base_url="http://localhost:11434")

# Universal invoke pattern
response = llm.invoke("Your prompt here")
result = llm.batch(["prompt1", "prompt2"])  # Parallel
async_result = await llm.ainvoke("prompt")   # Async
  \end{lstlisting}

  \docref{langchain-ai/langchain}
\end{frame}

% SLIDE 5: Prompt Templates
\begin{frame}[fragile]
  \frametitle{Prompt Templates — 3 Patterns}

  \begin{lstlisting}
# 1. Simple string template
from langchain_core.prompts import PromptTemplate
prompt = PromptTemplate.from_template("Explain {topic}")
prompt.invoke({"topic": "neural nets"})

# 2. Chat format (preferred for models)
from langchain_core.prompts import ChatPromptTemplate
chat_prompt = ChatPromptTemplate.from_messages([
    ("system", "You are a {role}"),
    ("human", "{question}"),
    MessagesPlaceholder(variable_name="history")
])

# 3. from_template shortcut
chat_prompt = ChatPromptTemplate.from_template(
    "System: {system}\nUser: {user}"
)
  \end{lstlisting}

  \textbf{Key:} Use \texttt{ChatPromptTemplate} for chat models, nest \texttt{MessagesPlaceholder} for history.

  \docref{Prompts}
\end{frame}

% SLIDE 6: LCEL Chains
\begin{frame}[fragile]
  \frametitle{LCEL Patterns (The 5 Most Common)}

  \begin{lstlisting}
# 1. Basic chain: prompt -> llm -> parser
chain = prompt | llm | StrOutputParser()
result = chain.invoke({"key": "value"})

# 2. Parallel branches (map)
from langchain_core.runnables import RunnableParallel
parallel = RunnableParallel(
    analysis=prompt | llm,
    summary=other_prompt | llm
)

# 3. Passthrough (preserve input)
from langchain_core.runnables import RunnablePassthrough
chain = ({"doc": RunnablePassthrough()} | llm)

# 4. Branching (conditional)
chain = condition | (true_branch | llm) or (false_branch | llm)

# 5. With fallbacks
chain = (prompt | llm).with_fallbacks([fallback_chain])
  \end{lstlisting}

  \docref{LCEL}
\end{frame}

% SLIDE 7: Tools & Agents
\begin{frame}[fragile]
  \frametitle{Tools \& Agents — Minimal}

  \begin{lstlisting}
# Define tool with decorator
from langchain_core.tools import tool

@tool
def multiply(a: int, b: int) -> int:
    """Multiply two numbers."""
    return a * b

# Bind tools to model
from langchain_core.tools import bind_tools
model_with_tools = llm.bind_tools([multiply])

# Create agent
from langchain.agents import create_react_agent, AgentExecutor
from langchain import hub
prompt = hub.pull("hwchase17/react")
agent = create_react_agent(llm, [multiply], prompt)
executor = AgentExecutor(agent=agent, tools=[multiply])

result = executor.invoke({"input": "What is 5 times 3?"})
  \end{lstlisting}

  \docref{Agents}
\end{frame}

% SLIDE 8: RAG Pipeline
\begin{frame}[fragile]
  \frametitle{RAG Pipeline — Complete Minimal Setup}

  \begin{lstlisting}
# 1. Load documents
from langchain_community.document_loaders import WebBaseLoader
loader = WebBaseLoader("https://example.com")
docs = loader.load()

# 2. Split
from langchain_text_splitters import RecursiveCharacterTextSplitter
splitter = RecursiveCharacterTextSplitter(chunk_size=1000)
splits = splitter.split_documents(docs)

# 3. Embed & store
from langchain_openai import OpenAIEmbeddings
from langchain_chroma import Chroma
embeddings = OpenAIEmbeddings()
vectorstore = Chroma.from_documents(splits, embeddings)

# 4. Retrieve & generate
retriever = vectorstore.as_retriever()
chain = {"context": retriever, "question": RunnablePassthrough()} \
    | prompt | llm | StrOutputParser()
answer = chain.invoke("What does the site say about X?")
  \end{lstlisting}

  \docref{RAG}
\end{frame}

% SLIDE 9: Memory & History
\begin{frame}[fragile]
  \frametitle{Memory Patterns}

  \begin{lstlisting}
# 1. RunnableWithMessageHistory (stateful chains)
from langchain_core.runnables.history import (
  RunnableWithMessageHistory
)
chain = RunnableWithMessageHistory(
    runnable=llm,
    get_session_history=get_session_history,
    input_messages_key="input",
    history_messages_key="history"
)
# Invoke with session_id
result = chain.invoke(
    {"input": "Hi!"},
    config={"configurable": {"session_id": "user123"}}
)

# 2. trim_messages (summarize old context)
from langchain_core.messages import trim_messages
trimmer = trim_messages(
    max_tokens=1000,
    strategy="last"
)

# 3. Store in any backend (Redis, PostgreSQL, etc.)
  \end{lstlisting}

  \docref{Memory}
\end{frame}

% SLIDE 10: LangGraph Essentials
\begin{frame}[fragile]
  \frametitle{LangGraph — State Machines}

  \begin{lstlisting}
from langgraph.graph import StateGraph
from typing import TypedDict

class AgentState(TypedDict):
    messages: list
    counter: int

# Build graph
graph = StateGraph(AgentState)

# Add nodes (functions)
def node_1(state):
    return {"messages": state["messages"] + ["step1"], "counter": 1}

def node_2(state):
    return {"messages": state["messages"] + ["step2"]}

graph.add_node("node1", node_1)
graph.add_node("node2", node_2)

# Add edges
graph.add_edge("node1", "node2")
graph.set_entry_point("node1")
graph.set_finish_point("node2")

# Compile & run
app = graph.compile()
result = app.invoke({"messages": [], "counter": 0})
  \end{lstlisting}

  \docref{LangGraph}
\end{frame}

% SLIDE 11: Streaming & Callbacks
\begin{frame}[fragile]
  \frametitle{Streaming \& Callbacks}

  \begin{lstlisting}
# Stream tokens
for chunk in llm.stream("Hello"):
    print(chunk.content, end="", flush=True)

# Async streaming
async for chunk in await llm.astream("Prompt"):
    print(chunk.content)

# Stream with events
from langchain_core.callbacks import LangChainTracer
from langsmith import Client
client = Client()
tracer = LangChainTracer(project_name="my_project")

result = chain.invoke(
    {"input": "..."},
    config={"callbacks": [tracer]}
)

# View in LangSmith
# https://smith.langchain.com -> projects -> my_project

# Custom callback
from langchain_core.callbacks import BaseCallbackHandler
class MyCallback(BaseCallbackHandler):
    def on_llm_end(self, response, **kwargs):
        print(f"LLM finished: {response}")
  \end{lstlisting}

  \docref{Callbacks, LangSmith}
\end{frame}

% SLIDE 12: Troubleshooting Checklist
\begin{frame}
  \frametitle{Top 10 Errors \& Fixes}

  \begin{enumerate}
    \item[\textbf{ImportError}] → Wrong module. Use \texttt{langchain\_openai}, not \texttt{langchain.llms}
    \item[\textbf{API Key missing}] → Set env var: \texttt{export OPENAI\_API\_KEY=...} or pass to constructor
    \item[\textbf{Chain invoke fails}] → Input dict keys must match prompt placeholders
    \item[\textbf{LLM returns None}] → Check \texttt{temperature}, \texttt{max\_tokens}, provider limits
    \item[\textbf{Parser error}] → Output format != parser expectation. Test LLM output directly.
    \item[\textbf{Tool binding fails}] → Tools must be List[BaseTool]. Use \texttt{@tool} decorator.
    \item[\textbf{Retriever returns empty}] → Vectorstore empty or wrong embedding model
    \item[\textbf{Memory not persisting}] → Check \texttt{session\_id} in config, backend connectivity
    \item[\textbf{LCEL pipe syntax error}] → Ensure left side outputs right side's expected input type
    \item[\textbf{LangGraph StateError}] → State keys must match TypedDict definition exactly
  \end{enumerate}

  \vspace{0.5cm}
  \textbf{Pro tip:} Always test components individually before chaining them.
\end{frame}

% SLIDE 13: Decision Matrix
\begin{frame}
  \frametitle{Decision Matrix: What to Use When}

  \begin{center}
    \small
    \begin{tabular}{l|c|c|c|c}
      \toprule
      \textbf{Scenario} & \textbf{Chain} & \textbf{Agent} & \textbf{LangGraph} & \textbf{Deep Agent} \\
      \midrule
      Fixed workflow & X & & & \\
      Dynamic tool use & & X & X & \\
      Multi-step reasoning & & X & X & X \\
      Parallel tasks & X & & X & X \\
      Complex state mgmt & & & X & X \\
      Human-in-loop & & & X & X \\
      Real-time streaming & X & X & X & \\
      Simple API call & X & & & \\
      RAG retrieval & X & X & X & \\
      Agentic loops & & X & X & X \\
      \bottomrule
    \end{tabular}
  \end{center}

  \vspace{0.5cm}
  \begin{tipbox}[Rule of Thumb]
    \small Start with \textbf{Chain}, add \textbf{Agent} for tool use, use \textbf{LangGraph} for fine-grained control.
  \end{tipbox}
\end{frame}

% SLIDE 14: Performance Tips
\begin{frame}
  \frametitle{Performance \& Best Practices}

  \begin{conceptbox}[Optimization]
    \begin{itemize}
      \item Use \texttt{.batch()} for parallel inference (faster than loop)
      \item Stream tokens with \texttt{.stream()} for UX (no waiting for full response)
      \item Cache embeddings—don't re-embed the same docs
      \item Use cheaper models for routing/filtering, expensive for final answer
      \item Set \texttt{temperature=0} for deterministic outputs (summaries, extraction)
      \item Chunk size 500–1500 tokens is sweet spot (not too small, not too large)
    \end{itemize}
  \end{conceptbox}

  \begin{warningbox}[Common Pitfalls]
    \begin{itemize}
      \item Don't reinitialize embeddings for each query (slow, inconsistent)
      \item Don't use retriever without persistence check (data loss)
      \item Don't forget to add error handling in agent loops (infinite loops)
    \end{itemize}
  \end{warningbox}
\end{frame}

% SLIDE 15: Resources & Next Steps
\begin{frame}
  \frametitle{Resources \& Next Steps}

  \begin{columns}[T]
    \column{0.5\textwidth}
    \textbf{Official Docs}
    \begin{itemize}
      \item \url{https://python.langchain.com}
      \item LangGraph: \url{https://langchain-ai.github.io/langgraph}
      \item LangSmith: \url{https://smith.langchain.com}
    \end{itemize}

    \textbf{Key Concepts}
    \begin{itemize}
      \item LCEL (Runnable chains)
      \item Output parsers
      \item Tool schemas
      \item State graphs
    \end{itemize}

    \column{0.5\textwidth}
    \textbf{Common Extensions}
    \begin{itemize}
      \item \texttt{langchain-openai}
      \item \texttt{langchain-anthropic}
      \item \texttt{langchain-community}
      \item \texttt{langgraph}
    \end{itemize}

    \textbf{Next: Deep Dive Into}
    \begin{itemize}
      \item Fine-tuning with function calling
      \item Multi-agent coordination
      \item Custom memory backends
      \item Production deployment
    \end{itemize}
  \end{columns}

  \vspace{1cm}
  \begin{center}
    {\large \textcolor{lcblue}{\textbf{This is a quick reference—bookmark it!}}}
  \end{center}
\end{frame}

\end{document}
